% !PRINTED 177
\weberpage{177}
The function $\psi(\omega)$ may possibly remain unchanged under certain
 transformations $S$ (e.\,g. always under the identity
 transformation). All transformations that leave a function
 $\psi(\omega)$ unchanged form a \gesperrt{group} $\mathfrak{G}$. For if
\[
\psi\,|\,S = \psi, \qquad \psi\,|\,S_1 = \psi,
\]
then also
\[
\psi\,|\,SS_1 = \psi\,|\,S_1 = \psi.
\]
The group $\mathfrak{G}$ is a subgroup of the group $\mathfrak{L}$ of all linear
transformations.

If $\mathfrak{G}$ is a subgroup of $\mathfrak{L}$ of \gesperrt{finite index $(\mathfrak{L}, \mathfrak{G})$, and
$\psi(\omega)$ a function of $\omega$ with the property that
conversely}
\[
\psi(\omega') = \psi(\omega)
\]
in turn entails \gesperrt{a linear connection}
\[
\omega' = \frac{\gamma+\delta\omega}{\alpha+\beta\omega}
\]
\gesperrt{in which $\begin{pmatrix} \alpha, & \beta \\ \gamma, & \delta \end{pmatrix}$ is a transformation belonging to $\mathfrak{G}$,
 then $\psi(\omega)$ is called a modular function belonging to the group $\mathfrak{G}$}.

We restrict ourselves to the consideration of those modular functions
that stand in an \gesperrt{algebraic connection} with the modulus $\varkappa^2$, and hence also with $j(\omega)$.

Under this hypothesis the following general theorem can be
stated:

\gesperrt{If $\psi(\omega)$ belongs to the group $\mathfrak{G}$ and $\chi(\omega)$ remains unchanged under the
transformations of $\mathfrak{G}$, then $\chi(\omega)$ is}
\textbf{rationally} \gesperrt{expressible through $\psi(\omega)$.}

For in the first place $\chi(\omega)$ is an algebraic function of $\psi(\omega)$,
and $\psi(\omega)$, as an algebraic function of $j(\omega)$, can take every
value. But to each value of $\psi(\omega)$ there belongs, by \gesperrt{hypothesis,
 only one value of $\chi(\omega)$, and therefore $\chi(\omega)$, as a one-valued
 algebraic function of $\psi(\omega)$, is rational.}

The study of the groups contained in $\mathfrak{L}$ and of the modular
functions belonging to them forms the principal subject of the great
work of \gesperrt{Klein and Fricke}: \glqq Vorlesungen über die Theorie
der elliptischen Modulfunktionen\grqq{} (2 Bde., Leipzig, Teubner,
1890, 1892). Here we consider only a few quite special cases
of these groups and functions, upon which we have come in the course of our
investigations.
